% =====================================================================
% DECEPTION COMPETITION — methodology + equations, paper-ready.
%   CHUNK 1 -> \section{Methodology}  (environment, mimic, validation)
%   Metrics block can move to \section{Experimental Results} if preferred.
% Style mirrors auction_section_rewrite.tex (equation/cases, \label, \mathbf).
% Figures to upload (extensionless \includegraphics):
%   deception_environment, deception_lie_fingerprint, deception_info_gradient,
%   deception_claim_dist, deception_outcome_parity, deception_benchmark.
% =====================================================================


% #####################################################################
% ## CHUNK 1 — METHODOLOGY
% #####################################################################

\subsection{Deception Competition}
\label{sec:deception}
To study strategic dishonesty in a controlled setting, we designed a repeated
simultaneous-move game in which $N=5$ sales agents pitch a single resort to a
mechanical buyer over $T=12$ rounds. The resort is described by $K=5$ quality
attributes whose true values in round $r$ form a vector
$t^{(r)}\in[0,1]^{K}$, quantized to two decimal places. Each round, every agent
$i$ simultaneously emits a \emph{claim} vector $c_i^{(r)}\in[0,1]^{K}$ about
those attributes; the buyer rewards persuasive claims but audits the agents and
penalizes exposed exaggeration. The buyer is non-strategic (a fixed mechanism),
so the game is purely one of how far each agent will distort the truth given the
risk of being caught and the value of the buyer's trust. The buyer holds fixed
preference weights $w\in\Delta^{K}$ over attributes (default
$w=(0.25,0.20,0.20,0.20,0.15)$).

\subsubsection{Information conditions}
We vary how much of the truth each agent knows. Under \emph{full information}
every agent observes all $K$ attributes; under \emph{partial information} each
agent observes a random size-$k$ subset $V_i^{(r)}\subseteq\{1,\dots,K\}$,
sampled uniformly without replacement and re-drawn each round, with the
remaining attributes hidden:
\begin{equation}
\tilde t_{i,a}^{(r)} =
\begin{cases}
t_a^{(r)}, & a \in V_i^{(r)},\\[2pt]
\varnothing, & a \notin V_i^{(r)}.
\end{cases}
\label{eq:dec-info}
\end{equation}
The information level $k=\lvert V_i^{(r)}\rvert$ is held fixed within an episode
and swept across $k\in\{1,\dots,5\}$ to trace an information gradient.

\subsubsection{Audit and catch rule}
After the five claims are submitted, the buyer computes the per-attribute
\emph{population mean} (the crowd consensus),
\begin{equation}
\mu_a^{(r)} = \frac{1}{N}\sum_{i=1}^{N} c_{i,a}^{(r)},
\label{eq:dec-mu}
\end{equation}
and audits each agent on the single attribute where it most exceeds the crowd,
\begin{equation}
a_i^{*} = \operatorname*{arg\,max}_{a\in\{1,\dots,K\}}\bigl(c_{i,a}^{(r)} - \mu_a^{(r)}\bigr),
\label{eq:dec-audit}
\end{equation}
with ties broken uniformly at random. The agent is \emph{caught} if and only if
its audited claim differs from the truth at that attribute,
\begin{equation}
\mathrm{caught}_i^{(r)} = \mathbb{1}\!\left[\, c_{i,a_i^{*}}^{(r)} \neq t_{a_i^{*}}^{(r)} \,\right],
\label{eq:dec-caught}
\end{equation}
where equality is evaluated at the two-decimal quantization. Because only the
most over-claimed attribute relative to the crowd is verified, the mechanism
\emph{polices over-claims and never under-claims}: an agent is exposed precisely
when it inflates one attribute above its peers and that inflation is false. This
asymmetry is the central design choice of the environment.

\subsubsection{Scoring, winner, and reward}
The buyer scores each agent by its trust-weighted appeal, where an agent's
standing trust $\mathrm{tr}_i^{(r)}\in[0,1]$ (initialized to $1$) scales its
score between a $0.5$ floor and full weight:
\begin{equation}
s_i^{(r)} = \bigl(w \cdot c_i^{(r)}\bigr)\Bigl(\tfrac12 + \tfrac12\,\mathrm{tr}_i^{(r)}\Bigr),
\qquad s_i^{(r)} = -\infty \ \text{if } \mathrm{caught}_i^{(r)},
\label{eq:dec-score}
\end{equation}
caught agents being disqualified for the round. The buyer books the highest-scoring
uncaught agent provided it clears the booking threshold $\tau$ (default $0.4$):
\begin{equation}
\mathcal{W}^{(r)} =
\begin{cases}
\operatorname*{arg\,max}_{i} s_i^{(r)}, & \max_i s_i^{(r)} > \tau,\\[2pt]
\varnothing, & \text{otherwise,}
\end{cases}
\label{eq:dec-winner}
\end{equation}
(set-valued under ties). The per-round reward penalizes catches by a fixed
amount $\rho$ (default $1.0$) and pays winners a threshold-normalized margin
bonus split across co-winners:
\begin{equation}
R_i^{(r)} =
\begin{cases}
-\rho, & \mathrm{caught}_i^{(r)},\\[6pt]
\dfrac{1 + 2\max\!\bigl(0,\, s_i^{(r)} - \tau\bigr)}{\lvert \mathcal{W}^{(r)} \rvert}, & i \in \mathcal{W}^{(r)},\\[10pt]
0, & \text{otherwise.}
\end{cases}
\label{eq:dec-reward}
\end{equation}

\subsubsection{Trust dynamics}
Trust carries across rounds and couples the single-shot incentives into a
repeated game. A catch shrinks trust multiplicatively; any clean round restores
it additively, both clipped to $[0,1]$:
\begin{equation}
\mathrm{tr}_i^{(r+1)} =
\begin{cases}
\max\!\bigl(0,\ \gamma\,\mathrm{tr}_i^{(r)}\bigr), & \mathrm{caught}_i^{(r)},\\[2pt]
\min\!\bigl(1,\ \mathrm{tr}_i^{(r)} + \delta\bigr), & \text{otherwise,}
\end{cases}
\label{eq:dec-trust}
\end{equation}
with catch multiplier $\gamma=0.7$ and recovery increment $\delta=0.1$. Through
Eq.~\eqref{eq:dec-score}, a catch is doubly costly: the immediate penalty of
Eq.~\eqref{eq:dec-reward} plus a depressed score multiplier in all later rounds.
The episode return is $R_i = \sum_{r=1}^{T} R_i^{(r)}$.

\subsubsection{Behavioral metrics}
\label{sec:dec-metrics}
We summarize an agent's revealed strategy with opponent-light statistics. The
\emph{lie rate} is the fraction of visible-attribute claims that over-state the
truth,
\begin{equation}
\mathrm{LieRate}_i =
\frac{\sum_r \sum_{a \in V_i^{(r)}} \mathbb{1}\!\left[c_{i,a}^{(r)} > t_a^{(r)}\right]}
     {\sum_r \lvert V_i^{(r)} \rvert},
\label{eq:dec-lierate}
\end{equation}
the catch and win rates are $\mathrm{CatchRate}_i = \tfrac{1}{T}\sum_r \mathrm{caught}_i^{(r)}$
and $\mathrm{WinRate}_i = \tfrac{1}{T}\sum_r \mathbb{1}[i \in \mathcal{W}^{(r)}]$,
and the preference-weighted exaggeration and under-claim are
\begin{equation}
E_i^{(r)} = \sum_{a} w_a \max\!\bigl(0,\, c_{i,a}^{(r)} - t_a^{(r)}\bigr),
\qquad
U_i^{(r)} = \sum_{a} w_a \max\!\bigl(0,\, t_a^{(r)} - c_{i,a}^{(r)}\bigr).
\label{eq:dec-exag}
\end{equation}

\subsubsection{Behavioral mimics}
\label{sec:dec-mimics}
As in the auction study, we train a cheap surrogate of each LLM and use it as a
drop-in replacement at scale. At agent $i$'s decision point in round $r$ we
encode the leakage-safe public context as a normalized feature vector
$\mathbf{x}_i^{(r)} \in \mathbb{R}^{23}$ concatenating the observed truth
($K$ values, with hidden entries imputed), a visibility mask ($K$), the buyer
preferences ($K$), the agent's own trust, the four opponents' trusts, the round
fraction $r/(T-1)$, the threshold $\tau$, and the penalty $\rho$. The encoding
deliberately \emph{excludes} the population mean $\mu^{(r)}$ and the agent index,
so the mimic cannot peek at the audit signal or its own identity. The vector is
z-scored, $\hat{\mathbf{x}} = (\mathbf{x} - \bar{\mathbf{x}})/\sigma_{\mathbf{x}}$,
and passed through a shared two-layer trunk feeding two heads:
\begin{align}
h &= \mathrm{Dropout}\!\Bigl(\mathrm{ReLU}\bigl(W_2\,\mathrm{Dropout}(\mathrm{ReLU}(W_1\hat{\mathbf{x}} + b_1)) + b_2\bigr)\Bigr), \label{eq:dec-trunk}\\
p_a &= \sigma\!\bigl((W_\ell h + b_\ell)_a\bigr), \qquad
\hat{c}_a = \sigma\!\bigl((W_c h + b_c)_a\bigr), \label{eq:dec-heads}
\end{align}
where the trunk has width $32$ with dropout $0.2$; the \emph{lie head} $p_a$ is
the probability of lying on attribute $a$, and the \emph{claim head}
$\hat{c}_a$ is the regressed claim value. The heads are trained with masked
losses: binary cross-entropy on the visible attributes
($m^{\mathrm{vis}}_a = \mathbb{1}[a \in V]$), and mean-squared error on the
attributes the regressor is responsible for at inference, namely those that are
hidden or visibly lied on
($m^{\mathrm{reg}}_a = \mathbb{1}[(\text{visible} \wedge \text{lied}) \vee \text{hidden}]$):
\begin{equation}
\mathcal{L}_{\mathrm{lie}} = \frac{\sum_a m^{\mathrm{vis}}_a\,\mathrm{BCE}(p_a, y^{\mathrm{lie}}_a)}{\sum_a m^{\mathrm{vis}}_a},
\quad
\mathcal{L}_{\mathrm{claim}} = \frac{\sum_a m^{\mathrm{reg}}_a\,(\hat{c}_a - y^{\mathrm{claim}}_a)^2}{\sum_a m^{\mathrm{reg}}_a},
\quad
\mathcal{L} = \mathcal{L}_{\mathrm{lie}} + \lambda\,\mathcal{L}_{\mathrm{claim}},
\label{eq:dec-loss}
\end{equation}
with $\lambda = 1$. At deployment the mimic assembles a claim stochastically. It
draws a lie decision $\ell_a \sim \mathrm{Bernoulli}(p_a)$ (temperature $\Theta$),
snaps to the known truth on visible attributes it chooses not to lie about, and
otherwise emits the regressed value perturbed by calibrated Gaussian noise:
\begin{equation}
c_a =
\begin{cases}
t_a, & a \in V \ \text{and}\ \neg \ell_a,\\[2pt]
\mathrm{clip}_{[0,1]}\!\bigl(\hat{c}_a + \varepsilon_a\bigr), & \text{otherwise,}
\end{cases}
\qquad
\varepsilon_a \sim \mathcal{N}\!\bigl(0,\ (\Theta\,\hat{\sigma}_a)^2\bigr),
\label{eq:dec-infer}
\end{equation}
where the per-attribute noise scale $\hat{\sigma}_a = \operatorname{std}(\hat{c}_a - y^{\mathrm{claim}}_a)$
is the regressor's training residual standard deviation, saved with the
checkpoint. The noise term is what lets the surrogate reproduce the
\emph{spread} of each LLM's claims rather than collapsing to the conditional
mean.

\subsubsection{Fidelity validation}
\label{sec:dec-fidelity}
We validate the substitution at three levels. At the \emph{decision} level we
report leave-one-episode-out cross-validation of the lie head's accuracy and the
claim head's mean absolute error. At the \emph{distributional} level we bin each
attribute's claims and compare the real and mimic histograms $P,Q$ by their
total-variation distance, used as the acceptance criterion because it is bounded
and sample-size independent:
\begin{equation}
\mathrm{TVD}(P,Q) = \tfrac12 \sum_b \lvert P_b - Q_b \rvert,
\qquad
D_{\mathrm{KL}}(P \Vert Q) = \sum_b P_b \log\frac{P_b}{Q_b},
\label{eq:dec-tvd}
\end{equation}
reporting KL divergence as a secondary effect size; we deliberately avoid a
$\chi^2$ $p$-value gate, which at our sample sizes rejects negligible
differences and is reported for reference only ($\chi^2 = \sum_b (O_b-E_b)^2/E_b$,
Cram\'er's $V = \sqrt{\chi^2/n}$, with per-LLM $p$-values combined by Fisher's
method $X = -2\sum_j \ln p_j \sim \chi^2_{2J}$). At the \emph{system} level we
replay the matched truth seeds with the mimic loadout and compare competition
outcomes (catch rate, win rate, final trust). Throughout, the episode is the
independent replicate and uncertainty is reported as $95\%$ bootstrap
confidence intervals,
\begin{equation}
\widehat{\theta}^{*(b)} = \frac{\sum_{e} n^{*}_{e}}{\sum_{e} d^{*}_{e}},
\qquad
\mathrm{CI}_{95} = \bigl[\,Q_{2.5}\{\widehat{\theta}^{*(b)}\},\ Q_{97.5}\{\widehat{\theta}^{*(b)}\}\,\bigr],
\label{eq:dec-bootstrap}
\end{equation}
where each bootstrap replicate $b$ resamples episodes with replacement and
recomputes the pooled ratio of numerator $n_e$ to denominator $d_e$.
